% 本文件由 LaTeX 排版 Agent 根据有效 translation.json 自动生成。
% author=Athanasius; work_id=2810; title=论狄奥尼修斯的观点

\chapter{论\Person{狄奥尼修斯}的观点}

% structure: p0001 source=h1 target=omitted reason=page-title-already-rendered-by-parent

\Emph{论\Person{狄奥尼修斯}的观点}\par

\Person{亚大纳削}关于\Person{狄奥尼修斯}——\Place{亚历山大}主教——的信函，表明他同\Place{尼西}大公会议一样反对亚略异端，并指出亚略派徒然诽谤他，声称他与他们立场一致。\par

1. \Emph{亚略派对\Person{狄奥尼修斯}的诉求是对他的诽谤}\par

你迟迟未告知我与基督的敌对着之间当前的争论；因为在你善意的信函到来之前，我已勤加探询，了解了此事，我闻之欣然。我赞许你的虔诚对我们有福的教父所持的正确意见，而在此次也再次认识到亚略派狂人的荒谬。因为他们的异端既无理性根据，亦无来自圣经的明确证据，他们便总是诉诸无耻的诡辩和似是而非的谬论。但他们如今竟敢诽谤教父：这并不奇怪，而是完全符合他们的邪恶。因为那些竟敢\BibleQuote{图谋反抗主，反抗他的受傅者}的人，现在又污蔑\Person{狄奥尼修斯}——\Place{亚历山大}的有福的主教——，声称他是他们的同党和共犯，这有什么可惊奇的呢？因为如果他们乐于抬举一个人，以支持自己的异端，即使他们称他为有福的，他们非但没有稍许贬损他，反而给他带来奇耻大辱；就如强盗或行为不端之人，当他们因自己的恶行而臭名昭著时，便宣称正派人士与他们同伙，从而污蔑正派人的清白。\par

2. \Emph{亚略派的立场与圣经相悖}\par

如果他们对己见的和主张有信心，就让他们赤裸裸地提出他们的异端，并说明他们认为从圣经或从人类理性中有什么宗教论证可为自己辩护。但如果他们拿不出任何东西，就让他们闭嘴。因为他们从任何方面都找不到辩护，只能使他们自己受到更大的定罪。首先，从圣经来看，因为\Person{若望}说：\BibleQuote{在起初已有圣言}，而他们却说：\Quote{他在受生之前并不存在}；而\Person{达味}以圣父的角色吟咏：\BibleQuote{我心灵洋溢美善的圣言}\BibleRef{咏 44:1}，他们却声称他只是思想中的存在，从无中而生。再者，\Person{若望}在\BibleBook{若望}{福音}中又说\BibleRef{若 1:3}：\BibleQuote{万物是藉着他而造成的；凡受造的，没有一样不是由他而造成的}，而\Person{保禄}写道：\BibleQuote{有一个主，就是耶稣基督，万物是藉他而有的}\BibleRef{格前 8:6}，及在别处：\BibleQuote{万物都是在他内受造的}\BibleRef{哥 1:16}，他们怎有胆量（不如说怎能免于羞辱）擅敢反对圣人们的言论，说万物的创造者是受造物，说那使一切受造物得以存在和存立的是一位受造物呢？其次，他们也没有任何从人类理性而来的宗教论证为自己辩护。因为有什么人，无论是希腊人还是蛮族人，会胆敢称他所宣认为天主的那位为受造物，或说他在被造之前并不存在呢？或者有什么人，当他听到他所相信为唯一天主的那位说：\BibleQuote{这是我的爱子}\BibleRef{玛 3:17}，及\BibleQuote{我心灵洋溢美善的圣言}时，竟敢说天主心中的圣言是从无中产生的呢？或说圣子是受造物，而不是说话者亲生的呢？再者，谁听到他所相信为主和救主的那位说：\BibleQuote{我在父内，父也在我内；我与父原为一体}\BibleRef{若 14:10}，竟敢将主所合一并保持不可分割的强行拆散呢？\par

3. \Emph{亚略派援引\Person{狄奥尼修斯}，就如犹太人援引\Person{亚巴郎}一样：但理由同样不堪}\par

因此，他们自己看到了这一点，并对自己的立场毫无信心，便对虔诚人士散布谎言。但对他们来说，更好的做法是，当他们孤立无援，并意识到在审查下他们理屈词穷、无言以对时，与其转身离开错误的道路，不去攀附他们并不认识的人，免得被那些人驳倒而蒙受更大的耻辱。然而，他们或许根本不想脱离自己的邪恶；因为他们效法了\Person{盖法}及其同伙的这一特点，正如他们从他们那里学会了否认基督。因为那些人也是如此：当主行了那许多事工，藉此显示自己是基督、是永生天主之子，并且被主定罪后，他们便开始事事处处违背圣经地思想和说话，片刻也无法面对针对自己的证据，便转向圣祖说：\BibleQuote{我们有亚巴郎为父}\BibleRef{玛 3:9}，妄图以此掩饰自己的荒谬。但他们这些话毫无收获，这些人也休想通过提及\Person{狄奥尼修斯}便能逃脱他人的罪责。因为主用这句话定了后者的恶行：\BibleQuote{亚巴郎却没有这样做过}\BibleRef{若 8:40}，而同样的真理也必将定这些人的不敬与虚假之罪。因为主教\Person{狄奥尼修斯}既不与亚略同流合污，也并非对真理无知。相反，无论是昔日的犹太人，还是当今的新犹太人，都从他们的父魔鬼那里继承了疯狂敌视基督的仇恨。那么，一个有力的证据再次表明这些人所言不实，是在诽谤此人，那就是：他既未因不敬而被其他主教定罪并逐出教会，如同这些人被开除圣职那样，也未作为异端分子自愿离开教会，而是在教会内光荣去世，他的纪念至今仍与教父们一同被保存和记录。因为如果他曾持有这些人的观点，或未证明自己所写内容为正确，那么毫无疑问，他也会遭受与这些人同样的对待。\par

4. \Emph{亚略派对\Person{狄奥尼修斯}的诉求基于他教导的一个孤立片段，而忽略其余}\par

事实上，这足以完全驳斥那些新犹太人，他们既否认主，又诽谤教父们，还试图欺骗所有基督徒。但既然他们认为，在这位主教的信函中某些部分，可以找到诽谤他的借口，来，让我们也看看这些部分，好让他们推理的虚妄即使是透过这些部分也暴露出来，使他们最终停止对主的亵渎，并且至少能与那些士兵\BibleRef{玛 27:54}一同，当他们目睹受造界作证时，宣认祂真是圣子，而非受造物之一。他们说，真福\Person{狄奥尼修斯}在一封信中曾说过：\Quote{圣子是一个受造物，是被造的，按本性不是祂自己的，而是在本质上与父相异的，正如农夫之于葡萄树，或造船者之于船一样，因为祂既是受造物，在祂存在之前并不存在。}是的，他写了这信，我们也承认他的信如此写道。但是，正如他写了这封，他也写了其他许多书信，他们应当也参考那些；为使这人的信仰能从所有书信中显明，而非仅凭这一封。因为一个建造许多三层桨战船的造船者的技艺，不是凭一艘船判断，而是凭所有。因此，如果他只是写了他们所说的这封信作为他信仰的说明，或如果这是他唯一的信，就让他们尽情指责他吧——因为这提议确实等于一项指控——但若他如此写是出于当时的情境与所涉之人，同时他也写了其他书信，针对受怀疑之处为自己辩护，那么他们就不该忽略这些理由，而草率地污蔑这人，免得他们显得只搜寻片言只语，却忽略了其他书信中所含的真理。因为农夫也根据所处理土壤的性质，以不同方式处理同类的树木：也不会有人因他砍伐一棵、嫁接另一棵、栽种一棵、再移动另一棵而责怪他。相反，一旦明了原因，他便更加钦佩其技艺的灵活多变。那么，除非他们肤浅地看了这著作，就让他们陈述这封信的主要主题吧；因为这样，他们计谋的恶毒与肆无忌惮便会显露出来。但既然他们不知道，或羞于陈述，我们必须自己来陈述。\par

\Emph{5. \Person{狄奥尼修斯}反对\Person{撒伯里乌}派写作的缘由。}\par

那时，在上利比亚的\Place{五城}地区，有些主教附从了\Person{撒伯里乌}。他们的主张极为成功，以致圣子几乎不再在教会中被宣讲。\Person{狄奥尼修斯}听闻此事，因他负责那些教会，便派人去劝诫那些犯错者停止他们的错谬，但他们非但不停止，反而在不虔敬上更加厚颜无耻，他被迫针对他们的无耻行径，写下了这封信，并从福音书中阐明救主的人性，以期既然那些人更加胆大地否认圣子，并将圣子的人性行动归于圣父，他便借此证明成为人的是圣子而非圣父，从而说服那些无知者圣父不是圣子，并逐步引他们认识圣子的真天主性与圣父的知识。这就是这封信的主要主题，也是他写此信的原因，正是因那些如此无耻地选择篡改真信仰之人。\par

\Emph{6. \Person{狄奥尼修斯}并未在所指摘的段落中表达其完整意见。}\par

那么，\Person{亚略}的异端与\Person{狄奥尼修斯}的主张之间有何共同之处呢？或者说，为何狄奥尼修斯要被称作与亚略相似，他们却大相径庭呢？因为一个是天主教会的导师，另一个则是一个新异端的创始人。当亚略为了阐述自己的错谬，以像埃及的\Person{索塔德}那样柔弱而可笑的风格写了一部\WorkTitle{塔利亚}时，狄奥尼修斯不仅也写了其他书信，还针对可疑之处为自己撰写了辩护，并清楚表明了自己持有正确意见。那么，如果他的著作前后矛盾，就让他们不要把他拉到自己那边，因为基于这个假设他就不值得信任。但若他在写给\Person{阿摩尼乌斯}的信中受到怀疑，便为自己辩护，以完善他先前所说的话，但这样做时并没有改变立场，那就显然可见，他是以一种受限定的意义来写那些受怀疑的段落的。然而，对于以如此意义所写或所行之事，人们无权恶意曲解，或各人按己意强解。因为即使是医生，也常根据其知识，对所处理的伤口施用某些在他人看来似乎不妥、却仅为健康着想的疗法。同样，一位明智的教师，其惯常做法是根据学生的性情来安排和讲授课程，直到把他们引上成全之路。\par

\Emph{7. 宗徒们的语言在特定段落中也需要类似的谨慎。}\par

但是，如果他们因为这位蒙福者这样写就指控他（因为\Person{亚略}派关于他的论证实际上就是对他的指控），那么当他们听到\WorkTitle{宗徒大事录}中伟大而蒙福的宗徒们的话时，又会怎样呢？先是\Person{伯多禄}说\BibleRef{宗 2:22}：\BibleQuote{诸位\Place{以色列}人，请听这些话：\Place{纳匝肋}人\Person{耶稣}，是天主以德能、奇迹和征兆——即天主藉祂在你们中所行的，正如你们自己知道的——向你们所举荐的人。祂照天主已定的计划和预知，被交付了；你们藉着不法者的手，钉祂在十字架上，杀死了祂。}又说\BibleRef{宗 4:10}：\BibleQuote{是凭\Place{纳匝肋}人\Person{耶稣}基督的名字，即你们所钉死，天主从死者中所复活的；就是凭着这人，这个站在你们面前的人好了。}以及\Person{保禄}在\Place{丕息狄雅的安提约基}叙述\BibleRef{宗 13:22}天主如何\BibleQuote{废了\Person{撒乌耳}，立了\Person{达味}为王，且为他作证说：\InnerQuote{我找到了\Person{叶瑟}的儿子\Person{达味}，他是一个合我心意的人，他要履行我的一切旨意。}天主按照恩许，从他的后裔中给\Place{以色列}兴起了一位救主\Person{耶稣}。}又在\Place{雅典}\BibleRef{宗 17:30}：\BibleQuote{天主对那愚昧无知的时代，原不深究；如今却传谕各处的人都要悔改，因为祂已定了一个日期，要由祂所立定的人，按正义审判天下，祂从死者中使他复活，给众人提供了信仰的保证。}或伟大的殉道者\Person{斯德望}说：\BibleQuote{看，我见天开了，并见人子站在天主右边。}啊，他们现在真是厚颜无耻（因为他们没有什么不敢做的），声称宗徒们本身持守\Person{亚略}的立场：因为他们宣称基督是来自\Place{纳匝肋}的人，且是能受苦的。\par

8. \Emph{宗徒们论及基督是人，但也是天主。}\par

那么，这些人既然有这种想象，难道宗徒们因为用了上述语言，就认为基督仅仅是一个人而非其他吗？绝对不可！这种想法根本不可能。然而，他们在这里也作为明智的建筑师和天主奥迹的管家行事。因为那时在错误中误导外邦人的犹太人，以为基督来临只是作为\Person{达味}后裔中的一个人，与\Person{达味}血统的其他子孙相似，既不愿相信祂是天主，也不信圣言成了血肉；因此，蒙福的宗徒们极有智慧，从向犹太人宣讲救主的人性特征开始，为的是通过对可见事实和所行奇迹的充分说服，使他们相信基督已经来临，然后引导他们上升到对祂天主性的信仰，指出祂所行的事不是出于人，而是出于天主。且看，\Person{伯多禄}称基督为能受苦的人，立刻在\BibleRef{宗 3:15}补充说：\BibleQuote{祂是生命之原，}而在福音中又宣认：\BibleQuote{你是基督，永生天主之子。}在他书信中称祂为\OriginalTerm{灵魂的主教}{Bishop of souls}，以及天使和掌权者并他自己的主。\Person{保禄}同样，称基督是\Person{达味}后裔中的人，在致\BibleBook{希伯来}{书}\BibleRef{希 1:3}中写道：\BibleQuote{祂是天主光荣的辉耀，及其本体的真像，}又在致\BibleBook{斐理伯}{书}\BibleRef{斐 2:6}中写道：\BibleQuote{祂虽具有天主的形体，并不以自己与天主同等作为应当把持不舍的。}然而，称祂为生命之原、天主子、辉耀、真像、与天主同等、主、灵魂的主教，这除了意味着在肉身内祂乃是天主的圣言，万物藉祂而造成，并且祂与圣父不可分离，如同辉耀与光不可分离一样，还能有什么别的意思呢？\par

9. \Emph{必须像宗徒们那样解释\Person{狄奥尼修斯}。}\par

\Person{狄奥尼修斯}因此依从他从宗徒们所学到的行事。因为\Person{撒伯里乌}的异端逐渐蔓延，正如我先前所说，他被迫写了上述书信，并向他们投掷救主在人性和谦卑方面所说的一切，为的是以祂的人性特征阻止他们说圣父是子，从而使他们更易接受关于圣子天主性的教导，而他在其他书信中从圣经称祂为圣言、智慧、能力、气息\BibleRef{智 7:25}，以及圣父的辉耀。例如，在他为自己辩护的书信中，他如我所述那样表达，在信德和对基督的孝爱中更显勇敢。因此，正如宗徒们不因为他们论及主的人性语言而受到指控——因为主已成为人——反而因其明智的保留和适时的教导而更受敬佩，同样，\Person{狄奥尼修斯}也不因其致\Person{欧夫拉诺}和\Person{亚摩尼乌}对抗\Person{撒伯里乌}的信而是亚略派。因为即使他使用了卑微的言辞和例证，但它们也出自福音，并且他使用它们的理由就是救主降生成人，因此不仅这些，还有许多类似的事被写了下来。因为正如祂是天主的圣言，所以后来\BibleQuote{圣言成了血肉}；在\BibleQuote{元始就是圣言}的同时，\BibleQuote{童贞女在世界穷尽时怀孕，主成了人}。而这两者所指示的是同一位，因为\BibleQuote{圣言成了血肉}。然而，论及祂的天主性和祂降生成人的表述，应当以分辨之心、切合具体语境地加以解释。那书写圣言人性特征的，也知道与祂天主性相关的事；那阐述祂天主性的，也并不忽略祂降生成人的事：但如同老练的和\OriginalTerm{被认可的钱币兑换者}{approved money-changer}，明辨每一者，从而走在纯正信仰的正道上；因此，当他说到祂哭泣时，他知道主既成为人，以流泪显示祂的人性，同时作为天主复活了\Person{拉匝禄}；知道祂肉体上饥渴，却以天主之能用五个饼喂饱五千人；知道一个人类的身体躺在坟墓中，却被圣言自己作为天主的身体复活起来。\par

10. \Emph{亚略派所引用的\Person{狄奥尼修斯}的表述是指基督作为人。}\par

狄奥尼修斯正是如此教导，在他致欧弗拉诺尔和阿摩尼乌斯的信中，针对撒伯里乌，论及救主的人性表述。因为属于后者的有这些话：\BibleQuote{我是葡萄树，我父是园丁}\BibleRef{若 15:1}，和\BibleQuote{向造祂的尽忠}\BibleRef{希 3:2}，以及\BibleQuote{祂造了我}\BibleRef{箴 8:22}，和\BibleQuote{成为比天使更尊贵的}\BibleRef{希 1:4}。然而他并非不知这些经文：\BibleQuote{我在父内，父在我内}\BibleRef{若 14:10}，以及\BibleQuote{看见我的，就是看见了父}。因为我们知道他在其他书信中提及了它们。当他在那里提及时，他也提了主的人性属性。因为就如\BibleQuote{祂虽具有天主的形像，却没有将与自己与天主同等视为不可攫取的宝藏，却空虚了自己，取了奴仆的形像}\BibleRef{斐 2:6}，以及\BibleQuote{祂本是富有的，却为了我们成了贫穷的}，因此，既有对祂天主性的崇高而丰富的描述，也有涉及祂降生成人的卑微而贫乏的表达。但下列理由明显表明，这些是就救主作为人而论的。园丁与葡萄树在性体上不同，而枝条与葡萄树是同一性体、同源的，实际上与葡萄树不可分割，葡萄树与枝条有同一个根源。但是，正如主所说，祂是葡萄树，我们是枝条。如果圣子与我们同一性体，且与我们有相同的根源，那么我们可以承认，在这一点上，圣子与圣父在性体上不同，就像葡萄树与园丁不同一样。但如果圣子与我们不同，祂是圣父的圣言，而我们是由土造成的，是亚当的后裔，那么上述表述就不应归于圣言的天主性，而应归于祂的人性降生。正因为如此，救主也说：\BibleQuote{我是葡萄树，你们是枝条，我父是园丁}。因为我们按肉身说与主是同族的，为此祂说：\BibleQuote{我要向我的弟兄宣告你的名}\BibleRef{希 2:12}。正如枝条与葡萄树同一性体，且出于葡萄树，我们也是这样，我们的身体与主的身体同质，从祂的满盈中我们领受了\BibleRef{若 1:16}，并以那身体作为我们复活和救恩的根。圣父被称为园丁，因为正是祂藉着自己的圣言栽培了葡萄树，即救主的人性，并藉着自己的圣言为我们预备了通往天国的道路；除非圣父吸引人归向祂，没有人能到主那里去\BibleRef{若 6:44}。\par

11. \Emph{宗徒们的类似语言也是如此。}\par

既然这表述的意义如此，那么对于这如此理解的葡萄树，确实经上写着：\BibleQuote{向那造了祂的尽忠}\BibleRef{希 3:2}，和\BibleQuote{成为比天使更尊贵的}\BibleRef{希 1:4}，以及\BibleQuote{祂造了我}\BibleRef{箴 8:22}。因为当祂取了那必须为我们奉献的，即取自童贞玛利亚的身体时，经上论到祂说祂被造、被形成、被制成：因为这类言辞适用于人。而且，并不是在取了肉身后祂才被造成比天使更尊贵的，以免显出祂先前不及或等同天使。但是，他写信给犹太人，将主的人性职分与梅瑟相比时，他说\BibleQuote{被造成比天使更尊贵的}，因为律法是藉天使颁布的，因为\BibleQuote{法律是藉梅瑟传授的，恩宠和真理却是由耶稣基督而来的}\BibleRef{若 1:17}，以及圣神的恩赐。并且，在那些日子，法律从丹宣讲到贝尔舍巴，现今\BibleQuote{他们的声音却传遍了大地}\BibleRef{罗 10:18}，外邦人朝拜基督，并藉着祂认识圣父。上述这些，都是就救主作为人而写的，并非其他方式。\par

12. \Emph{所引狄奥尼修斯的章节，若正确理解，是严格正统的。}\par

那么，难道如基督的仇敌反复声称的那样，狄奥尼修斯在论及圣子的人性特征，并因此称祂为受造物时，是意指祂是众人中的一个？或者当他说圣言并不符合圣父的性体时，是主张圣言与我们人同一性体？当然，他在其它书信中并非这样写，相反，在那些书信中，他不仅表明了正确的见解，而且还借这些书信竭力反对这些人，仿佛在说：我与你们这些天主的仇敌意见不同，我的著作也绝未为亚略提供不虔敬的借口。而是由于撒伯里乌派，我给阿摩和欧弗拉诺尔写信时，提及葡萄树和园丁，并使用其他类似的表达，为的是通过指出主的人性特征，来说服那些人不要说成为人的是圣父。因为正如园丁不是葡萄树，那在肉身内来的也不是圣父，而是圣言；圣言来到葡萄树内，因祂与枝条，即我们，有肉身上的亲缘关系，便被称为葡萄树。因此，我写给欧弗拉诺尔和阿摩尼乌斯的就是这个意思，但我以我写的其它书信来面对你们的厚颜无耻，好让心智健全的人知道这些书信所包含的辩护，以及我在基督信仰上的正确见解。因此，亚略派如果智力健全，就应该这样去思考和看待这位主教：\BibleQuote{因为凡明白的人，这一切都是显然的；凡寻得知识的人，这一切都是正直的}\BibleRef{箴 8:9}。但既然他们不理解天主教会的信仰，便陷入了不虔敬，结果，他们的智力受损，竟以为正直的事是弯曲的，称光明为黑暗，却以为黑暗是光明，因此，也有必要从狄奥尼修斯的其它书信中引述，并说明它们为何而写，以便更严厉地谴责异端者。因为正是从这些书信中，我们自己学会了去思考和撰写我们如今关于那位人物所作的论述。\par

13. \Emph{但狄奥尼修斯的其它著作也必须予以考虑。其历史。}\par

以下是导致他写其他信件的原因。\Person{狄奥尼修斯}主教得知\Place{五城}地区的情况后，出于对宗教的热忱，正如我在上文所说，写了信给\Person{欧福拉诺尔}和\Person{阿莫尼乌斯}，驳斥\Person{萨贝利乌斯}的异端。教会中有些弟兄，他们的意见是正确的，却没有事先询问他，以便亲自了解他信中的写法，就上\Place{罗马}去了；并且在与他同名的罗马主教\Person{狄奥尼修斯}面前说了他的坏话。罗马主教\Person{狄奥尼修斯}听到这些后，同时写信反对\Person{萨贝利乌斯}的追随者，也反对那些持有导致\Person{亚略}被逐出教会的观点的人；并称与\Person{萨贝利乌斯}或那些认为\OriginalTerm{天主圣言}{Word of God}是受造、被形成并有开端的人同流，是同等的、相反的不虔敬。他也写信给\Person{狄奥尼修斯}，告知他们对他所说的话。后者立刻写了回信，并给他的书题名为\WorkTitle{驳斥与辩护}。在此请注意可憎的基督仇敌之群，他们如何自招其辱。因为罗马主教\Person{狄奥尼修斯}也曾写信反对那些说\OriginalTerm{天主圣子}{Son of God}是受造物和被造之物的人，这表明并非从今日始，而是自古以来，\OriginalTerm{亚略派}{Arian}基督仇敌的异端就已被众人绝罚。而亚历山大的主教\Person{狄奥尼修斯}，就他所写的那封信进行辩护，表明他既不像他们所指责的那样思考，也从未持有过\OriginalTerm{亚略派}{Arian}的错误。\par

14. \Emph{狄奥尼修斯在\WorkTitle{驳斥与辩护}中的宗旨与一般方法}\par

\Person{狄奥尼修斯}仅就这些人之所纠缠的问题作出辩护，这一事实就足以完全定\OriginalTerm{亚略派}{Arians}的罪，并显示出他们的恶意。因为他写信不是出于愤怒的争论，而是要就自己受到怀疑的各点进行辩护。然而，他在针对指控进行辩护时，除了逐一驳斥每一项他所受到的怀疑，从而以这事实本身证实\OriginalTerm{亚略派}{Arian}狂人们的恶意之外，还做了什么呢？但为了通过他在辩护中所写的内容使他们完全混淆，来吧，让我将他的原话摆在你们面前。因为你们将从这些话中首先得知\OriginalTerm{亚略派}{Arians}是恶毒的，其次得知\Person{狄奥尼修斯}与他们的错误毫无关系。那么，首先，他将信写成了\WorkTitle{驳斥与辩护}的形式。但这无疑意味着，他旨在驳斥谬误的陈述，并就自己所写的内容进行辩护；表明他并非如\Person{亚略}所猜想的那样写信，而是在提及关于主在人性方面所说的事时，并非不知祂就是\OriginalTerm{圣言}{Word}与智慧，与\OriginalTerm{圣父}{Father}不可分离。然后他指责那些攻击他的人，因为他们没有完整地引用他的话语，而是断章取义，说话不诚实，而是心怀诡诈和任性。他将他们比作那些过去常挑剔真福宗徒书信的人。但他的这一控诉完全洗清了他身上的恶意嫌疑。因为如果他认为毁谤\Person{保禄}的人与毁谤他自己的人相似，他就恰恰表明，他是在\Person{保禄}的意义上写下那些话的。无论如何，在逐条回应对手的指控时，他解释了他们所引用的所有段落；并且，在后来的几封信中，他既在这些段落中推翻了\Person{萨贝利乌斯}，又表明他自己的信仰是多么纯正和虔诚。因此，尽管他们坚称\Person{狄奥尼修斯}持有这样的主张：\Quote{天主并非永远是圣父，圣子并非一直存在，而是天主没有圣言而存在，而圣子自己在受生之前并不存在：相反，曾有一个时刻祂不存在，因为祂不是永恒的，而是后来才获得存在}，——看他如何回答！他的话，无论是探究式的，还是综合推论式的，或是反问驳斥式的，或是对指控者的指责，由于他的论述冗长，我都略去了，只插入与对他的指控严格相关的内容。针对这些指控，他在题名为\WorkTitle{驳斥与辩护}的第一卷中，在若干前言之后，写了如下的话。\par

15. \Emph{选自\WorkTitle{驳斥与辩护}}\par

\Quote{因为从未曾有过一个时刻，天主不是圣父。}他在下文承认了这一点：\Quote{基督是永恒的，是\OriginalTerm{圣言}{Word}、智慧与大能。因为不可以为天主起初没有这样的子嗣，后来才生了圣子，而是圣子的存在并非源于自己，而是源于\OriginalTerm{圣父}{Father}。}稍后他又就同一主题补充道：\Quote{但作为永恒之光的光辉，祂自己当然也是永恒的；因为光既然常存，显然光辉也必常存。因为正是通过光照，光的存在才被察觉，没有不发光的光。让我们回到那些例子。如果有太阳，就有阳光，就有白日。如果没有这些，完全不可能有太阳。因此，如果太阳是永恒的，白日也将是无尽无休的。但事实上，并非如此，白日与太阳同始同终。然而，天主是永恒的光，无始无终。那么，光辉就永远在祂面前，与祂同在，无始而常生，在祂面前照耀，这正是那曾说：\BibleQuote{我原是祂的喜悦，日日在祂前欢跃，时时刻刻在祂面前欢跃}\BibleRef{箴 8:30} 的智慧。}稍后他又以这样的话回到同一主题：\Quote{圣父是永恒的，所以圣子也是永恒的，是光中之光：因为既有圣父，就必有圣子。如果没有圣子，如何能有圣父，又如何能指称为圣父呢？但两者都有，且是永恒地有。}然后他又补充道：\Quote{天主既是光，基督就是光辉；天主既是神，因为经上说：\BibleQuote{天主是神}\BibleRef{若 4:24}——同样，基督也被称为气息，因为祂是\BibleQuote{天主德能的气息}\BibleRef{智 7:25}。}再引第二卷的话，他说：\Quote{但惟独圣子，祂常与\OriginalTerm{圣父}{Father}同在，并充满于那\OriginalTerm{自有者}{IS}，祂自己也是从\OriginalTerm{圣父}{Father}而有的那\OriginalTerm{自有者}{IS}。}\par

16. \Emph{狄奥尼修斯用语与亚略用语的对比}\par

如果上述陈述的意思有疑义，就需要一位解释者了。但他既已清晰地、反复地论述同一主题，就让\Person{亚略}咬牙切齿吧，因为他看到自己的异端被\Person{狄奥尼修斯}推翻，听到他说出自己不愿听闻的话：\Quote{天主永远是圣父，圣子并非绝对永恒，而是祂的永恒由圣父的永恒所流溢，且与圣父同在，如同光辉与光同在。} 但是，那些竟敢想像\Person{狄奥尼修斯}与\Person{亚略}持同一立场的人，应当放下对他这样的诽谤。因为他们有什么共同之处呢？\Person{亚略}说，\Quote{圣子在受生之前并不存在，曾有一时，祂并不存在。} 而\Person{狄奥尼修斯}则教导，\Quote{天主是永恒的光，无始，亦永无终结：因此，光辉永恒地处于祂面前，与祂同在，在祂面前照耀，无始，且恒常受生。} 因为事实上，为了消除另一些人的猜疑——他们声称\Person{狄奥尼修斯}在论及圣父时不提圣子，又在论及圣子时不提圣父，而是将圣子与圣父分割、切除、分离——他就在第二卷中回应并羞辱了他们，如下文所述。\par

\Emph{\Person{狄奥尼修斯}并未分离天主圣三的位格。}\par

\Quote{我所提及的每一个名字，都与相邻的名字不可分离、不可分割。我提到了圣父，在提及圣子之前，我在圣父内也已暗指了祂。我提到了圣子——即使我没有明确提及圣父，圣父也必定在圣子内预先得到理解。我加上了圣神，但同时我也进一步加上了祂由谁及由何者所发出。但他们不知道，圣父\OriginalTerm{作为}{qua}圣父，并不与圣子分离——因为名字本身就带有这种关系——圣子也不会被逐离圣父。因为\InnerQuote{父}这个称呼本身就表示共同的联结。而在他们手中的是圣神，祂既不能与派遣祂者分离，也不能与传达祂者分离：那么，我既使用这些名字，又怎能想像是彼此分裂、完全隔绝的呢？} 然后过了一会儿他接着说，\Quote{因此，我们将\OriginalTerm{一体}{Monad}不可分割地伸展为\OriginalTerm{三位}{Triad}，又将\OriginalTerm{三位}{Triad}缩减而汇聚于\OriginalTerm{一体}{Monad}。}\par

\Emph{\Person{狄奥尼修斯}并未主张圣子与圣父不同性体。}\par

接下来，他针对他们指控他称圣子为受造物之一，且不与圣父同性同体（再次出现在第一卷中）而反驳他们，如下：\Quote{只是我在说某些事物被理解为受造的和被造的时，我粗略地举了一些不太恰当的例子，说植物与园丁、船与造船者并非同性同体。然后我更多地讨论了那些更贴切恰当的比较，并更详细地讨论了那些更接近真理的事物，提出了各种证据，这些我已经在另一封信中写给了你。通过这些证据，我也证明了他们对我的指控是不实的，即我否认基督与天主同性同体。因为虽然我指出我并没有找到这个词（\OriginalTerm{同性同体}{\GreekText{ὁμοούσιον}}），也没有在圣经的任何地方读到过它，但我随后提出的推理——这些推理被他们压制了——却与这个词的含义并不矛盾。因为我举了人类出生的例子，显然是为了表明同质性，并说明父母与子女的区别仅在于父母本身不是子女，否则就不存在父母或子女了。正如我之前所说，这封信由于情况不允许我出示，否则我本会把我当时所用的原话，或者说整封信的副本寄给你；如果有机会，我一定会这样做。但我记得，也回想起，我还加上了几个出自亲属关系的比喻。因为我说，一株植物由种子或根生出，虽然与它所从出的不同，但同时完全与其同性同体；一股溪流从泉源流出，获得了另一种形式和名称——泉源不叫河，河也不叫泉源——两者都存在，泉源如同父亲，而河则是泉源之水。但他们假装看不见这些以及类似的文字叙述，要像瞎子一样，却试图用两句不相干的话像石头一样从远处攻击我，不知道在超出我们知识、需要通过训练来理解的事情上，不仅外来的、甚至相反的例子也常常足以说明当前的问题。} 并且在第三卷中他说，\Quote{生命由生命而生，如河由泉源涌出，由不灭之光点燃了璀璨之光。}\par

\Emph{\Person{亚略}派诉诸\Person{狄奥尼修斯}之不连贯性。}\par

谁听了这些话，不把那些怀疑\Person{狄奥尼修斯}与\Person{亚略}持同一立场的人视为疯狂呢？看哪！在这些话中，他根据真理的论证践踏了整个\Person{亚略}异端。藉着\Quote{光耀}的比喻，他摧毁了\Quote{祂在被生之前不存在}和\Quote{曾有一时祂不存在}的说法，同样也藉着说圣父从未无子嗣。而他所谓的\Quote{从无中造}之主张，则被他以下的话摧毁：圣言如从泉源流出的河、如从根生出的枝、如父母所生的子女，是出自光明的光、出自生命的生命。而他们将圣言与天主隔开分离，则被他所说的\OriginalTerm{天主圣三}{Triad}毫无分割、毫无减损地聚敛于\OriginalTerm{单子}{Monad}内的话所推翻。至于他们说圣子与圣父的本质无分，他则毫不含糊地践踏，说圣子与圣父\OriginalTerm{同性同体}{of one essence}。在此，我们不得不惊异于这些不虔诚者的厚颜无耻。他们自称\Person{狄奥尼修斯}是他们同党，而\Person{狄奥尼修斯}却说圣子与圣父同性同体，他们怎能像蚊蝇般嗡嗡作响，抱怨\OriginalTerm{大公会议}{Synod}写下了\Quote{同性同体}是不对的呢？因为若\Person{狄奥尼修斯}是他们的朋友，他们就不得否认他们的同党所持守的。但若他们认为那措辞有误，又怎能反复坚称使用这措辞的\Person{狄奥尼修斯}与他们立场相同呢？更何况，他并非似乎只是顺便写下这些话，而是在先前已经写过其他书信，他揭穿那些指责他没有说圣子与圣父同性同体的人是撒谎；同时反驳那些认为他说圣言是受造的人，表明他并不持他们所猜想的主张，即便他曾使用过那些表述，也仅是为了显示是圣子而非圣父，穿上了那受造、形成、被造的身体；正因如此，圣子也被称为受造、被造和形成的。\par

20. \Emph{必须公正地解释\Person{狄奥尼修斯}，并允许他以自己的解释性陈述为自己辩护。}\par

显然，鉴于他先前已用过这些表述，同时彻底告别了\Person{亚略}派，他要求听者存有良心——他有资格辩称这些问题的难度，或者也许可以说其不可理解性——即他们不应评判文字，而应评判作者的意图，尤其是因为有许多证据可表明他的用意。例如，他亲自说：\Quote{我粗略地使用了这些关系的比喻，因为它们不太贴切，比如植物和园丁的比喻；而我详细论述了更恰当的比喻，并更详尽地探讨了更接近真理的比喻。}但是，说这话的人表明，称圣子为永恒且出于圣父，比称祂为受造，更接近真理。因为后者是指主的肉身本性，而前者是指祂\OriginalTerm{天主性}{Godhead}的永恒。例如，在接下来的话中，他坚称，且不仅坚称，而且刻意地、以真正的论证力量证明，那些指责他没有说圣子与圣父\OriginalTerm{同性同体}{of one essence}的人已被驳倒：\Quote{即使我在圣经中找不到这个措辞，但从圣经本身收集其一般意义后，我知道，作为圣子和圣言，祂不可能在圣父的\OriginalTerm{本质}{Essence}之外。}至于他不认为圣子是受造或形成之物——因为在这点上他们也反复引用他——他在第二卷中说：\Quote{但是，若我的诽谤者中有人因我称天主创造者为万有的制造者，就以为我意指祂也是基督的制造者，那么请他注意，我先前称祂为圣父，其中也暗含了圣子。因为当我说圣父是制造者后，我补充说，祂并不是祂所造之物的父，如果那位生者是严格意义上的父的话。因为\InnerQuote{父}这一术语的更广意义，我们将在后文阐述。圣父也不是制造者，如果制造者仅指工匠的话。因为在希腊人中，哲学家被称为自己言论的\InnerQuote{制造者}。而宗徒谈到\BibleQuote{法律的遵行者}\OriginalTerm{遵行者}{\GreekText{ποιητής}}\BibleRef{罗 2:13}，因为人们被称为内在品质的\InnerQuote{遵行者}，如美德和恶行；正如天主所说：\BibleQuote{我期望有人行正义，他却行邪恶}\BibleRef{依 5:7}。}\par

21. \Emph{\Person{狄奥尼修斯}在何种意义上说圣子是\Quote{被造成的}。}\par

确实，听到这些话的人会想起天主的神谕：\BibleQuote{恶人无论转向何处，必遭毁灭}\BibleRef{箴 12:7}。看哪！这些不虔诚的人巧妙地转向各个方向，终遭毁灭，甚至在涉及\Person{狄奥尼修斯}的事上也无借口可寻。因为他公开教导，圣子并非被造或受造之物，同时指责并纠正那些控告他说天主是（基督的）创造者的人，因为他们没有注意到他先前已称天主为圣父，其中也暗含了圣子。他这样说，就表明圣子并非受造物之一，天主并非祂圣言的制造者，而是圣父。而且，由于有人无知地反对他说他称天主为基督的制造者，他多方自我辩护，表明即便此处他所说的也无可指摘。因为他曾说，就基督的肉身而言，天主是制造者，因为那肉身是圣言所取，本身是受造的。但若有人猜想这是指圣言本身，他们也应给他公正的聆听机会。\Quote{因为我不认为圣言是受造物，也不称天主为祂的制造者，而称祂为圣父，即使我在提及圣子时顺便称天主为创造者，在此我仍能为自己辩护。因为希腊哲学家称自己为自己\OriginalTerm{言论}{\GreekText{λόγοι}}的\OriginalTerm{制造者}{\GreekText{ποιηταί}}，尽管他们是其父亲；而圣经描述我们甚至是我们内心活动的制造者（实行者），谈到法律的\InnerQuote{遵行者}以及审判与正义的\InnerQuote{遵行者}。}所以，他从各方面都表明，不仅圣子非被造或受造之物，而且他本人也与\Person{亚略}异端毫无瓜葛。\par

\Emph{22. 按照狄奥尼修，圣子与圣父的关系是本体的。}\par

因此，任何亚略派信徒都不应认为他讲过以下这类话：圣子与圣父共存，所以称谓虽然相关，但实体却相去甚远；并且因为圣子是后来才有的，圣子出生之后，天主由此获得了\Quote{父}这个附加名号，祂与圣子的共存也从此开始，如同在人间发生的情形那样。相反，他应当注意并牢记我们先前所讲的话，就会看到狄奥尼修的信仰是正确的。因为他说道：\Quote{因为从未有过天主不是圣父的时候}，又说：\Quote{无论如何，天主是永恒的光，无始无终，所以光辉永远在祂之前，与祂共存，无始而永受生，在祂面前照耀}——这些话应当使任何人都不可能对他抱有这类怀疑。此外，井与河、根与枝、气息与蒸气的比喻，也足以使基督的敌人在反复反对他时蒙羞。\par

\Emph{23. 狄奥尼修并不认为存在两位圣言。}\par

然而，亚略除了自己所有的不义，还从粪堆里扒拉出这种说法，补充说：\Quote{圣言并非圣父专有，而是在天主内的另一位圣言，而这一位——主——却在圣父的本体之外，与圣父的本体毫不相干，只是概念上被称为 \BibleQuote{圣言} ，并非按本性、按真理是天主子，而是被称为子，祂也如受造物一样，是出于收养义子。}——他这样说着，就在无知者中夸耀，好像在这方面也有狄奥尼修作他的同党；——那么请看看狄奥尼修关于这些问题的信仰，他是如何反驳亚略的这些邪说的。因为在第一册中，他这样写道：\Quote{我已说过，天主是一切善的泉源：而圣子被描绘成从祂流出的河。因为\InnerQuote{言语}是理智的流出，就借用于人身上的话说，那由舌头发出的理智，是从内心藉着口出来的，与心中的言语不同。因为后者发出前者之后，仍保持原样。而前者则被发出并飞去，向各处飘荡。因此二者互相寓居，又彼此不同：他们同是一，又同时是二。同样，圣父与圣子也被说成是一，彼此互相寓居。}在第四册中他又说：\Quote{正如我们的理智从自身发出言语，就像先知所说：\BibleQuote{我的心灵涌溢优雅的言辞}\BibleRef{咏 44:1}，并且二者彼此不同，各自占据不同的位置，一个居住在内心并在内心活动，另一个在舌上——然而它们并不分离，片刻也不失去彼此，理智并非没有表达（\OriginalTerm{无言}{\GreekText{ἄλογος}}），言语也并非没有理智，而是理智创造了言语，并在其中彰显自己，言语又显示理智原生其中，理智好似内在的言语，言语则是发出来的理智；理智过渡到言语，而言语将理智传递给听众：如此，理智藉着言语安居于听众的灵魂中，与言语一同进入；理智好似言语的父，自存自在，而言语则好似理智的子，其根源不在于理智之先，也不在于与理智同时从外而来，而是从理智涌出——同样，大能的圣父与普世理智，在万有之先就有圣子作为祂的圣言、诠释者与使者。}\par

\Emph{24. 如果亚略派信徒认同狄奥尼修，就让他们采用他的语言。}\par

这些事，亚略要么从未听过，要么听了却出于无知而不理解。否则，倘若他理解了，就不会如此严重地诽谤主教，而肯定会因他憎恨真理，也如咒骂我们一样咒骂他。因为他是基督的仇敌，必不犹豫地也迫害那些持守基督道理的人，正如主亲自预先说过的那样：\BibleQuote{若他们迫害了我，也要迫害你们}\BibleRef{若 15:20}。或者，如果那些不敬的首领们认为狄奥尼修是他们的同党，就让他们写下并宣认他所写的吧。让他们写下葡萄树与园丁、船与造船者的比喻；同时让他们也宣认，正如他在辩护中所做的那样，本体的合一，以及圣子出自圣父的本体，是永恒的；还有理智与言语、井与河等等的关系；好叫他们从鲜明的对比中看出，他使用前一类语言是有特殊用意的，而后一类才是表达基督信仰的完整含义。因此，让他们通过采纳后一种语言，撤销他们所持守的与之矛盾的观点。因为狄奥尼修的信仰怎会近似于亚略的毒害呢？亚略将\Quote{圣言}一词局限于概念意义，而狄奥尼修却按本性称祂为天主的真正圣言；一个将圣言从圣父驱逐出去，另一个却教导说圣言是圣父专有的，与祂的本体不可分离，如同言语之于理智，河之于井。那么，若有人能把言语与理智分离并驱逐，或把河与井分开，将它们隔开，或者说河的本质与井不同，并指出水是从别处来的，或胆敢将光辉与光分开，说光辉出自另一本质，那就让他去附和亚略的狂妄吧。因为这样的人连人的理智外表也将丧失。但是，如果本性知道这些是不可分的，而那些事物的产物就是它们自身所有，那么，就再没有人继续附和亚略或诽谤狄奥尼修了，反倒应基于这些理由，赞叹他语言的明白与信仰的正确。\par

\Emph{25. 狄奥尼修关于圣言的教导（续）。}\par

关于\Person{亚略}的疯狂，当他说在天主内的\OriginalTerm{圣言}{Logos}与\Person{若望}所说的那一位不同，即\BibleQuote{在起初已有圣言}\BibleRef{若 1:1}，并且天主在自己内的智慧与宗徒所说的\BibleQuote{基督是天主的德能和天主的智慧}\BibleRef{格前 1:24}不是同一个时，\Person{狄奥尼修}予以抵制并谴责此类错误，正如你在第二卷中所见，他在那里就这主题写道：\Quote{在起初已有圣言；但发出圣言的并非圣言，因为\BibleQuote{圣言与天主同在}。主成了智慧\BibleRef{格前 1:30}：那么发出智慧的那一位并非智慧，因为智慧说：\InnerQuote{我曾是他所喜悦的}。基督是真理：但他说：\InnerQuote{愿真理的天主受赞颂}\BibleRef{厄上 4:40}}。在此，他既推翻了\Person{撒伯里乌}，也推翻了\Person{亚略}，并指明这两种异端在不信中同等。因为圣言的父自身并非圣言，父的后裔也非受造物，而是由祂\OriginalTerm{本质}{ousia}所生的\OriginalTerm{独生子}{monogenes}。再者，发出的圣言不是父，也不是众多话语中的一言；而是独一的父之子，按本性是真实而真诚的儿子，祂现今在父内，并且永恒而无分离地出自父内。因此，主既是\OriginalTerm{智慧}{Sophia}又是真理，不在另一智慧之后居于第二位；而是唯有祂，父通过祂造了万有，并在祂内造了受造物的种种\OriginalTerm{本质}{ousia}，又通过祂将祂自己启示给祂所愿意的人，并在祂内运行和实施祂的普遍\OriginalTerm{眷顾}{pronoia}。\Person{狄奥尼修}只承认祂是天主的圣言。这就是\Person{狄奥尼修}的信德：因为我从他的书信中收集并抄录了几段陈述，足够促使你们增添数目，但为使亚略派因他们对主教的诽谤而彻底蒙羞。因为在他所写的一切，甚至细节中，他都揭露了他们的错误并痛斥了他们的异端。\par

26. \Emph{狄奥尼修如何对待撒伯里乌派}\par

因此，也很明显，即使写给\Person{优弗拉诺}和\Person{阿摩纽}的信，也是他出于不同意图和特殊目的而写的。因为他的辩护已清楚表明这一点。事实上，这是一种有效的论证方式，以推翻\Person{撒伯里乌}的疯狂，对于那些希望以简洁方式对付这些异端者的人，不应从适用于圣言天主性的表述入手，例如说子是天主的圣言、智慧和德能，以及\BibleQuote{我与父是一体}\BibleRef{若 10:30}，免得他们在听到这些经文时，歪曲这些正确的说法，作为他们厚颜无耻争辩的借口，当他们听到\BibleQuote{我与父是一体}，以及\BibleQuote{谁看见了我，就是看见了父}\BibleRef{若 10:30}时；而要强调论及救主作为人的所说的话，正如祂自己所做的那样，例如祂饥饿、口渴、困倦，以及祂是葡萄树，祂祈祷并受苦。因为就这些卑微的表述而言，这就更加显明，成为人的并非圣父。因为由此可推，当主被称为葡萄树时，必定也有一位园丁：当祂祈祷时，有一位俯听者；当祂祈求时，有一位给予者。如今，这样的事更迅速地显明\Person{撒伯里乌}派的疯狂，因为祈祷的那一位是一，俯听的那一位是另一，葡萄树是一，园丁是另一。因为论及子与父之间区别的无论什么表述，都是由于祂为我们承担的\OriginalTerm{肉身}{sarx}的缘故而用于祂。因为受造物在本性上与天主有别。因此，既然肉身是受造物，正如\Person{若望}所说，\BibleQuote{圣言成了血肉}\BibleRef{若 1:14}，尽管祂按本性是父所亲有的，与父不可分离，但由于肉身的缘故，父与祂有广大的分别。因为祂自己准许把适宜于肉身的事说在祂身上，为使显明这身体是祂自己的，而非其他任何人的。但这些话的含义就是如此，\Person{撒伯里乌}将会更快地被驳倒，因为这证明成为血肉的并非父，而是祂的圣言，这位圣言也赎回了肉身并将它奉献给父。如此驳倒并说服他之后，他就能更容易地教导他有关圣言的天主性，即祂是圣言和智慧，子和德能，光辉和\OriginalTerm{肖像}{charakter}。因为这里再次出现一个必然的推论：既然圣言存在，也就必然存在圣言的父；既然智慧存在，也就存在其根源；既然光辉存在，也就存在光明；这样，子和父也就是一。\par

27. \Emph{结论}\par

\Person{狄奥尼修斯}在写作时便知道这一点。他借着初期的书信使\Person{撒伯里乌}缄口无言，又在其它书信中战胜了\Person{亚略}的异端。因为正如救主的人性特质推翻了\Person{撒伯里乌}，同样，对抗亚略派的狂人，必须使用不是从人性特质，而是从彰显圣言天主性之处获取的证据，以免他们曲解因主的身体而论的言论，认为圣言与我们世人具有相同本性，从而继续执迷疯狂。但若他们也领受关于祂天主性的教导，他们就会谴责自己的错误；而一旦他们明白\BibleQuote{圣言成为血肉}，将来他们也就能更容易地区分人性特征与那些适合祂天主性的特征。既然事已如此，且\Person{狄奥尼修斯}主教因其著作显示出虔诚，亚略派的狂人还能做什么呢？在这证据前被定罪，他们还敢诽谤谁呢？因为他们既已从宗徒的根基上跌落，自己又没有定见，就必须寻找某种支持；若找不到，就诽谤教父们。但再没有人会相信他们，即使他们竭力加以中伤，因为该异端已受各方谴责。除非他们或许今后要谈论魔鬼，因为他是他们唯一的同伙，或者不如说，正是他将这异端灌输给他们的。那么，谁还能称那些以魔鬼为领袖的人为\Quote{基督徒}，而不称其为\Quote{\OriginalTerm{魔鬼徒众}{Diabolici}}，如此，他们所负的名号，不仅是基督的对敌，更是魔鬼的党徒呢？除非他们确实转变，摒弃自己所策划的不虔，认识真理。因为这对他们自身立刻便有益处，而且我们理当如此为所有陷于错谬的人祈祷。\par

\SourceRecord{Translated by Archibald Robertson.；Nicene and Post-Nicene Fathers, Second Series；Vol. 4.；Edited by Philip Schaff and Henry Wace.；Buffalo, NY: Christian Literature Publishing Co.,；1892.；<http://www.newadvent.org/fathers/2810.htm>.}
